% 按当前 GUI 源码核对的按钮与配置字段说明。
\newpage
\pagehead{9\quad 按钮速查：入口与主界面}{按界面英文原名查找。先分清配置、后台、硬件命令和运行记录。}
\subsection*{Run setup 与菜单入口}
\begin{tabularx}{\linewidth}{@{}L{63mm}Y@{}}\toprule
界面文字 & 含义、使用时机与实际影响 \\\midrule
\ui{Open YAML} & 读取 GetDAQ 连接与运行设置（YAML），填入节点、IP、端口、路径和实验名。它不是打开板卡 XML 的按钮。\\
\ui{Save YAML} & 把当前 GetDAQ 设置写入所选 YAML。保存设置不会将 XML/XCFG 发送给 ECC。\\
\ui{Nodes / Paths} & 打开节点配置表、源配置目录、GUI 记录目录及超时设置；详见第 11 页。\\
\ui{Master / Slave} & 选择本机角色并进入服务配置。单机也选 Master，然后 Use this computer。只改变角色选择不代表后台已经启动。\\
\ui{Configuration} & New configuration… 复用服务另存新 YAML，Run 归零并清空记录参数。User guide 查看中英说明。\\
\ui{ECC / data service instances…} & 打开同一套主从与 ECC 服务管理界面；可以连接已有 ECC，也可以创建后台组。\\
\ui{GET configuration files...} & 启动独立的原版配置编辑器，编辑 Configuration source 目录中的 XML/XCFG。当前界面要求先 Disconnect。\\
\ui{Tools} & 工具入口：配置同步、直连 CoBo 信息、在线/离线波形、文件转换和日志导出。\\
\ui{ECC / data service manager…} & 与配置菜单的 ECC / data service instances 是同一管理入口，不是再启动一套后台。\\
\ui{Run} & Pause run 在实际 Running 时暂停；Resume run 在实际 Paused 时继续。重连外部采集也可操作；恢复延续原 Run。\\
\ui{Recovery → Status and warnings…} & 查看并确认警告、检查当前状态，或明确执行退回 Idle；详见第 19 页。\\
\bottomrule\end{tabularx}
\subsection*{主界面输入与记录参数}
\begin{tabularx}{\linewidth}{@{}L{63mm}Y@{}}\toprule
界面文字 & 含义、使用时机与实际影响 \\\midrule
\ui{Experiment / Run} & 实验名和本次 Run 编号。实验名参与记录及归档命名；Describe/Prepare/Configure 显式 ID 优先决定 ECC 配置。\\
\ui{Run description / note} & 本次采集备注，随运行记录保存。\\
\ui{Add parameter} & 在 Run record parameters 中增加一行实验记录参数，不配置板卡寄存器。\\
\ui{Remove selected} & 删除参数表中当前选中的记录参数行。Nodes 中的同名按钮作用于节点行，见第 11 页。\\
\bottomrule\end{tabularx}
\note{最常用入口：已有后台 → Master → Connect existing ECC；软件管理后台 → Master → Create ECC / data services。从机先选 Slave 并启用。}

\newpage
\pagehead{10\quad 按钮速查：采集与自动运行}{这些按钮会操作硬件或本次 Run；不要和后台进程的 Start / Stop 混淆。}
\subsection*{Acquisition controls}
\begin{tabularx}{\linewidth}{@{}L{63mm}Y@{}}\toprule
界面文字 & 含义、使用时机与实际影响 \\\midrule
\ui{Connect} & 按当前节点表连接 ECC，开始查询真实状态与接收监控；受管配置可先检查后台。连接本身不执行 Describe / Prepare / Configure。\\
\ui{Disconnect} & 断开 GUI 的控制/观察会话。仍有工作或采集时按提示确认；远端后台与硬件状态保留。\\
\ui{Progress all} & 对所有节点推进一个初始化阶段，依次 Idle → Described → Prepared → Ready；重复到 Ready。会发送对应硬件命令。\\
\ui{Regress all} & 对所有节点退回一个初始化阶段，执行相应 Undo / Breakup。不是关闭后台服务。\\
\ui{Progress selected node} & 只推进下方选择框中的一个节点；先选对节点。\\
\ui{Regress selected node} & 只让选择框中的一个节点退回一个初始化阶段。\\
\ui{Run start} & 开始本次采集并建立运行记录；若已观察到 Running，则附着观察。已有同号记录或接收文件时按确认提示处理。\\
\ui{Run stop} & 请求停止采集；ECC 确认 Ready 后处理归档、备份和本地记录。服务进程仍运行。\\
\ui{Retry finalization} & 重试可恢复的本地记录收尾；不等于再发送 Stop，也不直接重放结果不确定的 catalog。\\
\ui{Recover journal} & 读取已有运行日志 journal，恢复中断后的记录处理；不会凭空认定硬件已经停止。\\
\ui{Resume failed catalog replacement} & 仅特定同号替换失败、仍在原会话且检查通过时出现。继续原替换流程；普通归档警告不使用它。\\
\bottomrule\end{tabularx}
\subsection*{Automatic Run}
\begin{tabularx}{\linewidth}{@{}L{63mm}Y@{}}\toprule
界面文字 & 含义、使用时机与实际影响 \\\midrule
\ui{Enable / Time / File size} & Enable 启用规则；Time 设置时长阈值；File size 设置接收数据量阈值。两项都选时，先到阈值者触发。\\
\ui{Stop / Next Run} & 阈值达到后，Stop 只停止本次；Next Run 在成功收尾、确认 Ready 后进入下一次。\\
\ui{Runs / Unlimited} & 自动序列总次数，包含第一次；0 显示 Unlimited。\\
\ui{Apply} & 在 Run start 前应用自动运行规则。仅修改输入还没有应用；手动 Run stop 取消自动续跑。\\
\bottomrule\end{tabularx}
\note{日常顺序：Connect → Progress 至 Ready → Run start → Run stop。需要关闭 ECC / 接收进程，再进入服务管理使用 Stop，而不是只 Disconnect。}

\newpage
\pagehead{11\quad 配置按钮：Nodes / Paths}{这里保存 GUI 连接设置；板卡的电子学参数仍在原版 XML/XCFG 中。}
\subsection*{按钮与路径}
\begin{tabularx}{\linewidth}{@{}L{63mm}Y@{}}\toprule
界面文字 & 含义、使用时机与实际影响 \\\midrule
\ui{Configure and manage ECC / data services…} & 进入主从角色、已有 ECC 连接、后台创建/管理；接受配置后带回 Run Setup。\\
\ui{Add CoBo} & 在当前节点表新增一行，填写该节点对应 ECC 及接收器。不会启动一个 ECC 进程。\\
\ui{Remove selected} & 移除节点表当前行；不删除远端实例、配置文件或原始数据。\\
\ui{Browse…} & 选择它左边输入框的本机目录。配置源与 GUI 记录目录各有一个，不是远端目录浏览器。\\
\ui{Sync configuration to project folders before starting services} & 保存启动时的同步开关。勾选后由托管启动流程分发配置；仅勾选并保存 Nodes 不立即传文件。\\
\ui{Save and apply} & 检查节点/路径并保存到当前 YAML；新设置会创建自己的 YAML。应用到 GUI，不直接同步 XML 或配置硬件。\\
\ui{Cancel} & 放弃本窗口尚未保存的表格与路径修改。此前已接受的远端后台操作仍保留。\\
\bottomrule\end{tabularx}
\subsection*{节点表：每一列填什么}
\begin{tabularx}{\linewidth}{@{}L{63mm}Y@{}}\toprule
界面文字 & 含义、使用时机与实际影响 \\\midrule
\ui{ID / Role / Name} & GUI 唯一节点编号、cobo 或 mutant、显示名称。托管多 ECC 可各有 CoBo[0]，其实际板卡索引另行保留。\\
\ui{ECC host / ECC port} & 运行 SOAP ECC 的电脑地址及 SOAP 服务端口，例如 8083。这里不填 CoBo 的板卡 IP。\\
\ui{Sentry host / Sentry port} & 运行接收目录监控及归档服务的电脑地址、HTTP 端口，例如 8080。\\
\ui{Router host / Router port} & 板卡发送原始数据的接收地址及数据端口，例如数据网 IP:46005；不是接收器控制端口。\\
\ui{Legacy ECC target host / port} & 旧 ECC 可选第二条板卡数据目标。正常单入口配置端口留 0；这不是 DataExporter 的下游导出端口。\\
\ui{Describe / Prepare / Configure ID} & ECC 用于匹配配置的 ID。例：cobo0 → describe-cobo0；pedestals → configure-pedestals。\\
\ui{Preview host / Preview TCP port} & 独立在线波形的 DataExporter TCP 导出地址，常见 46491；0 表示不配置预览。\\
\bottomrule\end{tabularx}
\note{Configuration source folder 是本机 XML/XCFG 源目录；GUI Run records root 保存 YAML/journal/CSV/备份；Raw .graw folder 是 Sentry 返回的接收主机目录。三者用途不同。}

\newpage
\pagehead{12\quad 配置按钮：连接已有 ECC}{Master → Connect existing ECC；用于已有的本机或其他电脑上的 SOAP ECC。}
\subsection*{选择、检查与保存}
\begin{tabularx}{\linewidth}{@{}L{63mm}Y@{}}\toprule
界面文字 & 含义、使用时机与实际影响 \\\midrule
\ui{Saved ECC connections / New connection} & 选择已保存的连接，或新填一套。选择条目只载入设置，不证明服务在线。\\
\ui{Remove saved connection} & 从历史选择列表移除该条目。保留对应 YAML、远端服务和数据。\\
\ui{Existing receiver IP and data port} & 展开已有接收器、在线预览和 GUI 记录目录设置；不是启动接收器的开关。\\
\ui{Browse…} & 选择本机 Local Run records directory，保存 GUI 运行记录。\\
\ui{Check connection} & 探测 SOAP 协议并读取 ECC 状态；不会初始化板卡，也不会启动 ECC 或接收器。修改地址等输入后要重新检查。\\
\ui{Save and use in Run Setup} & 检查成功后保存连接 YAML 并带回主界面；当前入口会安排连接查询。\\
\ui{Stop waiting} & 结束本地检查等待，远端 ECC 服务保持原样。\\
\ui{Close (keep services running)} & 关闭窗口，已运行的后台继续运行。\\
\bottomrule\end{tabularx}
\subsection*{输入框含义}
\begin{tabularx}{\linewidth}{@{}L{63mm}Y@{}}\toprule
界面文字 & 含义、使用时机与实际影响 \\\midrule
\ui{Instance name} & 便于识别的连接名称，不是板卡配置文件 ID。\\
\ui{ECC computer IP / ECC service port} & 已有 ECC 所在电脑地址和 SOAP 监听端口。不同实例通过主机地址与端口区分。\\
\ui{Data receiver IP / port (SOAP ECC)} & 填写板卡应发送到的接收数据地址/端口。后台真实监听应事先设置一致。\\
\ui{Receiver type} & 选择 DataRouter / no TCP preview 或 DataExporter。这里只描述已有接收器，不切换远端程序。\\
\ui{TCP preview export port} & 仅 DataExporter 在线波形使用的导出端口，与板卡接收端口分开。\\
\ui{Local Run records directory} & 当前 GUI 电脑的运行记录路径；不是远端原始数据输出目录。\\
\bottomrule\end{tabularx}
\note{外部 ECC 仍读取它自己 --config-repo-url 指向的目录或 NFS。这个连接页不会把本机 XML 自动发送到未托管的外部 ECC。多个已有 ECC 的连接参数可在 Nodes 表统一查看和编辑。}

\newpage
\pagehead{13\quad 配置按钮：主控、主机与配置源}{Master → Create ECC / data services；先确定目录与主机，再创建板卡组。}
\subsection*{主控页面与高级配置服务}
\begin{tabularx}{\linewidth}{@{}L{63mm}Y@{}}\toprule
界面文字 & 含义、使用时机与实际影响 \\\midrule
\ui{Configuration source folder / Browse…} & 选择主控电脑上保存完整 XML/XCFG 及附属资源的目录。此目录是分发源，不是从机上的路径。\\
\ui{Sync configuration to project folders before starting services} & 启动托管后台前同步源目录到执行主机项目目录。关闭时使用执行主机已有文件，并显示未核验警告。\\
\ui{Sync project configuration…} & 立即选择本机/从机同步完整源配置；不要求启动采集。详见第 17 页。\\
\ui{First-time setup / edit hosts and boards} & 展开 Hosts、ECC and boards 与 Advanced: master service；默认勾选。隐藏这些页不会停服务。\\
\ui{Start / check master service} & 启动或检查主控的配置分发服务。不是启动 SOAP ECC，也不是硬件 Start。\\
\ui{Open existing project…} & 选择含 project.json 的主控项目目录，恢复主机、组和部署记录；在线状态还需 Refresh。\\
\ui{Stop project configuration service} & 停止主控配置分发服务；ECC、接收器和节点管理服务保留。\\
\bottomrule\end{tabularx}
\subsection*{Hosts 页}
\begin{tabularx}{\linewidth}{@{}L{63mm}Y@{}}\toprule
界面文字 & 含义、使用时机与实际影响 \\\midrule
\ui{Project data folder on this computer / Browse…} & 启用本机执行节点前选定工作目录。ECC、接收器和 Sentry 在此工作；已启用主机保留原目录。\\
\ui{Use this computer (automatic)} & 把主控自己作为执行主机，建立/检查本机节点管理与项目目录。单机部署先点它。\\
\ui{Add slave} & 粘贴从机 Copy pairing code 生成的完整邀请码，再加入该执行主机；不是填写 CoBo 的 IP。\\
\ui{Name} & 主机显示名称。主机表还显示管理地址、身份、项目数据目录及最后观察状态。\\
\bottomrule\end{tabularx}
\subsection*{Advanced: master service 输入}
\begin{tabularx}{\linewidth}{@{}L{63mm}Y@{}}\toprule
界面文字 & 含义、使用时机与实际影响 \\\midrule
\ui{Master project directory} & 主控保存项目描述、版本、凭据和部署回执的目录，不是原始数据目录。\\
\ui{Master IP reachable from other hosts} & 从机能够访问的主控电脑地址，用于 HTTPS 配置服务。\\
\ui{Configuration service port} & 主控配置分发端口，常见 24250；不是 SOAP ECC 端口。\\
\bottomrule\end{tabularx}
\note{主控源目录 → 配置版本 → 每个执行主机缓存 → 各自的 Project data folder。启用后的工作目录不会因为修改 GUI 记录目录而自动迁移。}

\newpage
\pagehead{14\quad 配置按钮：ECC 组与启动范围}{ECC and boards 表支持 Ctrl 多选、Shift 连选。组代表后台实例及其板卡/接收器配置。}
\subsection*{日常启动与采用配置}
\begin{tabularx}{\linewidth}{@{}L{63mm}Y@{}}\toprule
界面文字 & 含义、使用时机与实际影响 \\\midrule
\ui{ECC configuration} & 上方选择框选一个组；选择它会切换表格到该组。要选多个组，在 ECC and boards 表中多选。\\
\ui{Start project services} & 使用当前表格选择，重新读源配置，准备并启动所选组及对应接收器、共享 Sentry；然后带回 Run Setup。\\
\ui{Start selected groups and use} & 与上面的启动采用相同流程；明确按多选范围操作。不同组通过独立实例/端口区分，不启动未选组。\\
\ui{Use this ECC in Run Setup} & 采用当前选中、已准备且活动的组；可多选。它不负责启动尚未运行的组。\\
\ui{Refresh / check original operation} & 查询主机、进程及已接受操作的回执。历史 failed 是历史记录，不等于当前服务仍离线。\\
\ui{Stop waiting} & 结束本地等待，不撤销远端已接受的准备、启动或停止。之后用 Refresh 查询。\\
\bottomrule\end{tabularx}
\subsection*{组的增删与停止}
\begin{tabularx}{\linewidth}{@{}L{63mm}Y@{}}\toprule
界面文字 & 含义、使用时机与实际影响 \\\midrule
\ui{Add boards / ECC…} & 新增组草稿，填写主机、端口、目录、实验 ID、CoBo index 和接收器类型。按 OK 保存草稿，不立即改变远端服务。\\
\ui{Edit selected group…} & 编辑当前组的草稿；该按钮针对当前一组，不批量编辑所有选中行。\\
\ui{Remove selected group from next release} & 从下一版本项目草稿移除当前组；活动后台需先正常停止或核查释放。不会删除原始数据与审计。\\
\ui{Save and sync only} & 保存并准备项目版本，分发配置但不激活后台。这一按钮处理项目中所有组，不采用表格的多选启动范围。\\
\ui{Stop selected group} & 正常停止当前一组后台，并非所有多选组。先在主界面结束采集，再关闭后台。\\
\ui{Terminate selected group…} & 确认后终止当前组匹配的后台进程。只用于正常停止无法完成时；进程退出不代表板卡已停。\\
\ui{Review terminated group…} & 确认板卡停止、数据保留后，节点再核验进程确实不存在，释放旧部署占用；不删除旧数据和回执。\\
\ui{Stop project services} & 正常停止项目中活动的受管后台组。采集中的组应先 Run stop。\\
\ui{Close (keep services running)} & 关闭管理窗口，后台独立继续运行。\\
\bottomrule\end{tabularx}
\note{多选启动后，各组共同进入一个 Run Setup，应使用相同 Experiment。启动后台只准备进程与配置；Describe / Prepare / Configure / Run start 仍由采集流程发出。}

\newpage
\pagehead{15\quad 配置字段：新增或编辑 ECC 组}{对应 Add boards / ECC… / Edit selected group… 中的每一项。}
\subsection*{身份、路径与执行主机}
\begin{tabularx}{\linewidth}{@{}L{63mm}Y@{}}\toprule
界面文字 & 含义、使用时机与实际影响 \\\midrule
\ui{Instance name} & 自定义组名，便于从多个 ECC/接收器中识别。\\
\ui{Configuration source (master; edit above)} & 主控源配置目录；本对话框只读，要到上层 Configuration source 修改。\\
\ui{Experiment / configuration ID} & 如 pedestals；用于选择原版配置 ID 和运行命名。应存在 configure-pedestals.xml 或 .xcfg。\\
\ui{ECC execution host} & 选择执行 SOAP ECC 的电脑；来自已加入的本机/从机列表。\\
\ui{Project folder on ECC host} & 该 ECC 所在电脑的工作目录，收到的完整配置会同步到这里。\\
\ui{Data execution host} & 选择执行 DataExporter / DataRouter 的电脑，可与 ECC 相同或不同。\\
\ui{Project folder on receiver host} & 接收器工作目录，原始数据写入这里；归档后通常位于该目录下的实验/run 子目录。\\
\ui{Receiver} & 选择 DataExporter 或 dataRouter；只有实际提供 TCP 导出功能的接收器可用于在线预览。\\
\bottomrule\end{tabularx}
\subsection*{地址、端口与板卡索引}
\begin{tabularx}{\linewidth}{@{}L{63mm}Y@{}}\toprule
界面文字 & 含义、使用时机与实际影响 \\\midrule
\ui{ECC computer IP / SOAP ECC port} & ECC 监听的电脑控制网地址与 SOAP 端口，常见 8083。\\
\ui{Alarm callback port} & 原版 ECC 的报警回调端口；同机多实例也要避免冲突。\\
\ui{Receiver control IP / control port} & 接收器电脑的控制网地址与管理端口，常见 46003。\\
\ui{Receiver data IP} & CoBo 能到达的接收电脑数据网地址；可以与控制 IP 不同。\\
\ui{CoBo data destination port} & 板卡数据接收端口，常见 46005。应与 ECC 数据目标一致。\\
\ui{Waveform TCP export port} & DataExporter 转发给独立波形客户端的下游端口，常见 46491。选择 dataRouter 时不使用此项。\\
\ui{Sentry monitor port} & 监控该项目目录的 HTTP 端口，常见 8080。同一主机同一项目目录共享 Sentry。\\
\ui{CoBo index} & 该 ECC 对应的实际 CoBo[N] 索引；多个 ECC 可以各自控制 CoBo[0]。GUI 节点编号可能不同。\\
\ui{OK / Cancel} & OK 校验并保存组草稿；Cancel 放弃本次编辑。保存草稿后，还需 Start 或 Sync 才会改变执行主机配置。\\
\bottomrule\end{tabularx}
\note{端口属于不同进程和用途。同机多实例必须规划不冲突的 ECC、报警、接收控制/数据、导出端口。不同电脑可以使用相同端口号。}

\newpage
\pagehead{16\quad 配置按钮：从控 Slave}{从机只需先选本机目录并启用；实验配置由主控整目录分发。}
\subsection*{按钮和勾选项}
\begin{tabularx}{\linewidth}{@{}L{63mm}Y@{}}\toprule
界面文字 & 含义、使用时机与实际影响 \\\midrule
\ui{Slave: receive configuration and run backend services} & 把当前电脑作为执行节点，接收配置并运行后台；不要求在从机再手工选择主控的 XML 文件。\\
\ui{Browse…} & 选左侧对应的本机路径：缓存目录或项目数据目录。缓存选择器会在选定目录下使用 getdaq-node-cache。\\
\ui{Enable / check slave} & 保存/检查本机节点配置并启动管理服务。启用节点不等于已启动 ECC、接收器或采集。\\
\ui{Copy pairing code} & 复制尚未消费且在有效期内的邀请码，到主控 Add slave 粘贴；配对后此按钮不可用是正常状态。\\
\ui{Renew pairing code} & 尚未配对或邀请码过期时生成新码。部分旧标签显示 Renew / re-pair…。\\
\ui{Remove master connection…} & 已配对时显示。确认后撤销旧主控管理凭据并生成新码；已接受工作先完成，后台和数据保留。\\
\ui{Management network settings (optional)} & 展开节点管理地址与端口；通常使用默认值，只在需要指定可达网卡时更改。\\
\ui{Remove selected inactive entry} & 隐藏实例列表中选中的无进程历史项；只是从本机列表移除，不删除节点审计/数据。仍有实际进程的项会保持可见。\\
\ui{Close (keep services running)} & 关闭窗口。节点管理、ECC、接收器可继续独立运行。\\
\bottomrule\end{tabularx}
\subsection*{目录、地址和状态表}
\begin{tabularx}{\linewidth}{@{}L{63mm}Y@{}}\toprule
界面文字 & 含义、使用时机与实际影响 \\\midrule
\ui{Configuration cache (on this slave)} & 软件缓存、配对、部署版本及回执的本机目录，不是主控的路径。\\
\ui{Project data folder (on this slave)} & 该从机后台工作及原始数据目录。收到配置后自动同步到此目录；已启用的主机保留设置。\\
\ui{Local management IP / port} & 主控可访问的从机管理地址/HTTPS 端口，常见 24260；不是 ECC SOAP 或板卡数据端口。\\
\ui{Instance / Type / Process} & 后台实例身份、ecc/receiver 类型，以及是否发现匹配进程。Process found 不等于板卡已 Running。\\
\ui{Local record / ECC configuration cache / Raw data} & 部署记录阶段、ECC 使用目录、原始数据工作目录；可与主控路径摘要逐项核对。\\
\bottomrule\end{tabularx}
\note{从机无需多开主界面来运行多个 ECC。主控为它创建多个后台组即可。主机配对与硬件连接是两层：配对成功还需启动服务，再通过 Run Setup 连接 ECC。}

\newpage
\pagehead{17\quad 配置同步、超时与常用路径}{保存 YAML、同步 XML、让板卡加载参数是三个不同操作。}
\subsection*{同步对话框与实际影响}
\begin{tabularx}{\linewidth}{@{}L{63mm}Y@{}}\toprule
界面文字 & 含义、使用时机与实际影响 \\\midrule
\ui{Tools → Sync project configuration…} & 重新读取主控源目录，选择执行主机。主控自己的执行节点也可以作为目标。\\
\ui{主机复选框} & 选择接收本次配置的主机，可多选；只有目标主机本地的项目目录会被处理。\\
\ui{Sync selected hosts} & 上传不可变版本，执行主机下载/校验后复制到本机项目数据路径。保留资源相对路径，并备份被替换的配置。\\
\ui{OK / Cancel} & 对话框确认按钮显示 Sync selected hosts；Cancel 放弃此次目标选择，不撤销以前已完成的同步。\\
\ui{活动实例更新确认} & 接受警告后可以在稳定 Ready / Running 下同步文件；已知命令忙时暂缓替换。同步不把硬件强制退回 Idle。\\
\bottomrule\end{tabularx}
\subsection*{Nodes / Paths 中的时间设置（单位 ms）}
\begin{tabularx}{\linewidth}{@{}L{63mm}Y@{}}\toprule
界面文字 & 含义、使用时机与实际影响 \\\midrule
\ui{Status polling} & 周期状态查询间隔；数值小会更频繁地访问 ECC/Sentry。\\
\ui{HTTP timeout} & 普通 HTTP 请求等待上限。此项不是硬件状态转换总时限。\\
\ui{Catalog timeout (server normally waits 30 s)} & 等待 Sentry 归档响应的上限；服务器归档自身可能等待约 30 秒。\\
\ui{Transition deadline} & 一次硬件转换及状态确认的等待时限；超时后可能远端已经执行，应先查状态。\\
\ui{Post-stop settling time} & 硬件停止后等待接收数据落稳的间隔，再进入收尾处理。\\
\bottomrule\end{tabularx}
\subsection*{四个目录分别放什么}
\begin{tabularx}{\linewidth}{@{}L{63mm}Y@{}}\toprule
界面文字 & 含义、使用时机与实际影响 \\\midrule
\ui{Configuration source folder (master)} & 你用配置编辑器修改的完整实验配置；这是复制源。\\
\ui{Project data folder} & 执行主机工作目录：ECC 配置副本、接收原始 .graw、Sentry 归档实验/run 子目录。\\
\ui{Configuration cache} & 软件管理缓存、版本和部署回执；收到文件后再同步到工作目录。\\
\ui{GUI Run records root} & 当前 GUI 的 run.yml、journal、CSV 与配置备份；与原始数据目录独立。\\
\bottomrule\end{tabularx}
\note{修改 Nodes 并 Save and apply 不会立即传 XML。修改 XML 后手动 Sync，或在勾选同步时 Start project services。文件传完后，板卡仍需对应 Configure 命令才会加载新参数。外部未托管 ECC 继续使用原 NFS/仓库。}

\newpage
\pagehead{18\quad 原版配置编辑器：按钮与菜单}{Configuration → GET configuration files...；此窗口编辑源文件，保存后再同步。}
\subsection*{文件、编辑和显示}
\begin{tabularx}{\linewidth}{@{}L{63mm}Y@{}}\toprule
界面文字 & 含义、使用时机与实际影响 \\\midrule
\ui{New / Open} & 新建配置或打开配置文件。New 是新文件，不是新建一次采集。\\
\ui{Save / Save As... / Save All} & 保存当前文件、另存为新路径、保存全部已打开文件。保存到主控 Configuration source 才会进入后续项目分发。\\
\ui{Quit / 标签页关闭} & 退出编辑器或关闭当前文档；有未保存修改时按提示处理。\\
\ui{Undo / Redo} & 撤销/重做编辑。EditConfigDialog 中也有对应图标按钮；它们不会撤销已发到硬件的命令。\\
\ui{Find... / Find Next / Search...} & 打开属性搜索或查找下一个匹配项。Search 也可从配置树右键进入。\\
\ui{Case sensitive / Whole words} & 搜索区分大小写/要求完整匹配项；具体匹配按原编辑器规则。\\
\ui{Keep open} & 找到属性后保留搜索窗口，便于继续查找。\\
\ui{Name / Value / Unit / Range} & View 中切换属性名称、值、单位、范围列的显示；只改显示，不改采样参数。\\
\ui{Backtrack level / Include URL} & 显示继承回溯层级与外部引用地址，帮助判断参数从哪个文件继承。\\
\ui{Attribute builder / Insert} & 打开属性构造面板，按输入的名称、类型和值插入属性。编辑对话框的 builder 图标也用于打开该面板。\\
\ui{About} & 查看原版编辑器信息。\\
\bottomrule\end{tabularx}
\subsection*{配置树右键菜单}
\begin{tabularx}{\linewidth}{@{}L{63mm}Y@{}}\toprule
界面文字 & 含义、使用时机与实际影响 \\\midrule
\ui{New... / Override / Rename...} & 新增节点/属性、在当前配置覆盖继承项、重命名当前项。Override 会改变当前层的有效参数。\\
\ui{Delete / Copy / Cut / Paste} & 删除、复制、剪切、粘贴配置树内容；删除有实际配置影响，保存前核对。\\
\ui{Expand all / Collapse all} & 展开/折叠配置树，不改变文件参数。\\
\ui{Import / Export → File...} & 将文件内容导入配置或把选定配置导出到文件，注意与 Open/Save 的文档范围区分。\\
\ui{Include → File...} & 添加文件引用，保留依赖关系；被引用资源也应放在完整源目录内。\\
\ui{Database...} & 原版保留的数据库入口，是否可用取决于已构建的存储插件；普通 XML 文件流程无需它。\\
\ui{Inheritance 子菜单} & 查看 Overrides / Overriden by / Inherits / Inherited by / Defined in 等继承来源与关联项。\\
\bottomrule\end{tabularx}
\par{\footnotesize 编辑后：Save → 主控 Sync → ECC Configure；外部未托管 ECC 继续使用原仓库。}

\newpage
\pagehead{19\quad 工具按钮：波形、日志与状态}{在线查看独立于 GUI；清屏与接受警告不代表硬件状态变化。}
\subsection*{离线 / 在线波形查看器}
\begin{tabularx}{\linewidth}{@{}L{63mm}Y@{}}\toprule
界面文字 & 含义、使用时机与实际影响 \\\midrule
\ui{Open waveform... / Previous / Next} & 离线打开 .graw/.gmat；Previous/Next 切换帧。Event ID + Go to event 查事件，Frame 选择其帧，AGET 选择芯片。\\
\ui{Auto scan / Period} & 离线按 Period 周期自动扫描帧；取消勾选停止扫描。\\
\ui{DataExporter / Refresh choices} & 在线选择当前配置中的一个或全部导出器，自动使用导出地址；Refresh choices 重新读 GUI 提供的列表。\\
\ui{Pause preview / Follow latest} & Pause preview 暂停在线预览连接；Follow latest 跟随最新缓存事件，取消后可查看固定事件。\\
\ui{Event ID / Go / Source / CoBo / AsAd / AGET} & 在线从缓存选择或跳转事件，再筛选来源与电子学通道。未缓存的旧事件应读离线数据。\\
\ui{Channels… / All / None / OK / Cancel} & 勾选显示通道，All 全选，None 全不选，OK 应用，Cancel 放弃。\\
\ui{3D / Align baselines / Export plot/data…} & 3D 切换纵深图；Align baselines 默认减去每通道中位数并按零参考等距显示，幅度为 $\Delta$ADC，取消则用绝对 ADC。PNG 保存画面，CSV 保留原始 ADC。\\
\ui{New session} & 在线清除预览缓存并开始新预览会话；原始数据与采集 Run 保留。\\
\bottomrule\end{tabularx}
\subsection*{日志与警告}
\begin{tabularx}{\linewidth}{@{}L{63mm}Y@{}}\toprule
界面文字 & 含义、使用时机与实际影响 \\\midrule
\ui{Event log / ECC log / Data log} & 切换当前日志类别；数据日志按实例显示 DataExporter/DataRouter，右侧选择主机/实例或共享 Sentry。\\
\ui{Clear / Export combined logs…} & Clear 只清当前显示。Tools → Logs 导出本次 GUI 会话收集的 TXT，按接收时间与序号合并，并标明来源。\\
\ui{Acknowledge selected} & 确认看过选中通知。不会证明不确定的硬件命令已完成。\\
\ui{Accept selected warning} & 接受选中的一般警告并记住处置；失败或结果不确定的动作仍需核查。\\
\ui{Clear acknowledged notifications / Clear displayed log} & 清除已确认的通知/事件框显示，保留持久日志和控制器状态。\\
\ui{Recheck state} & 更新本窗口的 ECC 阶段显示，依据 GUI 当前观察状态；不是硬件复位按钮。\\
\ui{Reset to Idle} & 确认后对已停止、符合条件的 ECC 逐步发送 Undo/Breakup，退回 Idle。与接受警告不同，会操作硬件。\\
\ui{Close} & 关闭状态窗口，不改变硬件与后台。\\
\bottomrule\end{tabularx}
\note{每次 GUI 启动使用新日志会话。清屏后同一会话的导出历史仍保留。远端日志原文时间保留，合并顺序采用 GUI 接收时间，不能当作完全同步的跨主机时钟。}

\newpage
\pagehead{20\quad 工具按钮：文件、CoBo 与操作示例}{补充工具不会替代 ECC 采集流程；下面给出可照着操作的入口顺序。}
\subsection*{Data file tools…}
\begin{tabularx}{\linewidth}{@{}L{63mm}Y@{}}\toprule
界面文字 & 含义、使用时机与实际影响 \\\midrule
\ui{模式选择框} & Inspect raw frames 检查原始帧；Convert to waveform matrix 生成 .gmat；Merge events 生成 .gmer。FITS/ROOT 类选项依赖相应已安装工具。\\
\ui{Add files… / Remove selected} & 加入已完成采集的本机输入文件，或从本次工具输入列表移除所选项。移除列表项不删原文件。\\
\ui{Browse…} & 选择该工具的输出路径。转换使用新的输出位置，先核对模式和输入。\\
\ui{Run tool / Stop tool} & 运行所选独立数据工具/停止该工具子进程；不控制 ECC 采集。\\
\ui{Close} & 关闭工具窗口；工具仍工作时按窗口提示处理。\\
\bottomrule\end{tabularx}
\subsection*{CoBo information (direct, read-only)…}
\begin{tabularx}{\linewidth}{@{}L{63mm}Y@{}}\toprule
界面文字 & 含义、使用时机与实际影响 \\\midrule
\ui{Connect / refresh CoBo} & 输入板卡的硬件控制地址和端口后，直连查询 CoBo 信息；这里填 CoBo IP，不填 SOAP ECC 电脑 IP。\\
\ui{Disconnect observation} & 断开该只读信息工具的观察连接；不执行 Run stop，也不关闭 ECC。\\
\ui{Close} & 关闭信息窗口。该工具不执行硬件配置或采集命令。\\
\bottomrule\end{tabularx}
\subsection*{三个最常用操作}
\begin{tabularx}{\linewidth}{@{}L{63mm}Y@{}}\toprule
界面文字 & 含义、使用时机与实际影响 \\\midrule
\ui{单机：启动软件管理的后台} & Master → Create ECC / data services → 选择源目录 → Hosts / Use this computer → Add boards / ECC → 选组 → Start selected groups and use → 主界面 Progress 至 Ready → Run start。\\
\ui{多机：让从机运行后台} & 从机 Slave 选本机目录并 Enable → Copy pairing code → 主控 Add slave → 新建组选择该从机 → 主控多选组启动。配置会先传到从机缓存，再放入从机项目目录。\\
\ui{修改电子学配置} & 配置编辑器 Open/修改/Save 到源目录 → Tools / Sync project configuration → 勾选目标主机 → Sync selected hosts → 需要时通过 ECC Configure 加载。\\
\bottomrule\end{tabularx}
\note{通用弹窗：OK/Yes 接受本次明确动作；Cancel/No 放弃本次选择。Stop waiting 只停止本地等待，不能据此判断远端操作被取消。灰色按钮应结合窗口状态提示判断，先 Refresh/查看日志，不反复发同一硬件命令。}
